\chapter{How Refusal Was Reproduced}
\label{ch:reproduced}
Everything that follows is built. This chapter is about what was never built and reproduced itself anyway, which is the book's evidence that it can be.

Every instrument in the rest of the book has a drafter: the floor, the review ratio, the schools, the charter, the record of a bearer's birth. The people who refused in chapter~\ref{sec:1}'s record had none. Browning's policemen, Milgram's subjects, the Oliners' rescuers, the maroons and the people who ran the Underground Railroad were not produced by any institution designed to produce refusers, and most of them refused against the institutions that governed them. So the record can be asked a plain question: what did they have, at each of the four stages, and who supplied it?

The stages are the Introduction's. Something makes refusing survivable, or the whole price falls on whoever refuses. Somebody forms the refuser before the occasion, because nobody in the room when it comes can put the refusal there. Something keeps the refuser going between one decision and the next, because a refusal made once and then worn down is only a delay. And somebody hands the refusal to people who weren't there, or each generation starts from nothing. The rest of the book builds them in that order, walls first. The record supplied them in nearly the reverse order. For most of the people in it the last three stages came first, from their own people, and the first came, when it came at all, from somebody else's order.

Two cautions about the evidence. The record holds the refusals that left a record: a policeman's statement to a postwar investigator, a rescuer still alive forty years later to be interviewed, a burial society whose dues and dead somebody counted. The refusers who were killed before anyone wrote them down are missing from it, and so are the societies that were found out. The record therefore overstates how often refusal survived, and it can't say how often refusal was tried. And the evidence is mixed in kind: two laboratory rooms, a police battalion, a survey of rescuers, a Red Cross delegate's account of one village, histories of movement schools and of a sharecroppers' union, a compilation of Black economic history Du Bois published in 1907. None of it was gathered to answer my question, and where a source supports less than the claim I put on it, I say so.

There are two lineages in it. The operator's is the officer's: people inside an apparatus, asked to do its work, some of whom said no. The acted-upon's is the lineage of the people the apparatus was built to sort. The Introduction said the older of the two reproduced itself with no institution the order built. This chapter is where that claim gets tested, one stage at a time, and where it has to be qualified. What the book takes from the reading is a list of what each stage needed, read off people who had it. A machine starts with nothing on the list. The record can't show that the list can be supplied by design. It shows what a design would have to supply.
